% Edge Detection in Digital Image Processing
% Topics: Gradient operators, Canny edge detection, Laplacian methods, performance metrics
% Style: Technical report with computational implementations

\documentclass[a4paper, 11pt]{article}
\usepackage[utf8]{inputenc}
\usepackage[T1]{fontenc}
\usepackage{amsmath, amssymb}
\usepackage{graphicx}
\usepackage{siunitx}
\usepackage{booktabs}
\usepackage{subcaption}
\usepackage[makestderr]{pythontex}
\usepackage{algorithm}
\usepackage{algpseudocode}

% Theorem environments
\newtheorem{definition}{Definition}[section]
\newtheorem{theorem}{Theorem}[section]
\newtheorem{example}{Example}[section]
\newtheorem{remark}{Remark}[section]

\title{Edge Detection: From Gradient Operators to Multi-Scale Analysis}
\author{Computer Vision Laboratory}
\date{\today}

\begin{document}
\maketitle

\begin{abstract}
This report presents a comprehensive analysis of edge detection algorithms in digital image
processing. We examine gradient-based methods (Sobel, Prewitt, Roberts), the Canny edge
detector with its multi-stage pipeline, and Laplacian-based approaches including zero-crossing
detection. Computational implementations demonstrate edge detection on synthetic test patterns
and natural images, with quantitative performance evaluation using precision, recall, and
F-measure metrics. We analyze the trade-offs between noise sensitivity, localization accuracy,
and computational efficiency across different detector families.
\end{abstract}

\section{Introduction}

Edge detection is a fundamental operation in computer vision, providing critical information
about object boundaries, surface discontinuities, and scene structure. An edge represents a
significant local change in image intensity, typically occurring at the boundary between
regions of different brightness or texture.

\begin{definition}[Edge]
An edge in a digital image is a set of connected pixels that form a boundary between two
regions with relatively distinct intensity properties. Mathematically, edges correspond to
local maxima in the intensity gradient magnitude:
\begin{equation}
|\nabla I| = \sqrt{\left(\frac{\partial I}{\partial x}\right)^2 + \left(\frac{\partial I}{\partial y}\right)^2}
\end{equation}
where $I(x,y)$ is the image intensity function.
\end{definition}

\section{Theoretical Framework}

\subsection{Gradient-Based Edge Detection}

The image gradient captures the rate and direction of intensity change at each pixel.

\begin{definition}[Image Gradient]
For a continuous image $I(x,y)$, the gradient is a vector:
\begin{equation}
\nabla I = \left[\frac{\partial I}{\partial x}, \frac{\partial I}{\partial y}\right]^T
\end{equation}
with magnitude $|\nabla I| = \sqrt{I_x^2 + I_y^2}$ and direction $\theta = \arctan(I_y/I_x)$.
\end{definition}

\begin{theorem}[Sobel Operator]
The Sobel operator approximates the gradient using 3$\times$3 convolution kernels:
\begin{equation}
G_x = \begin{bmatrix} -1 & 0 & 1 \\ -2 & 0 & 2 \\ -1 & 0 & 1 \end{bmatrix}, \quad
G_y = \begin{bmatrix} -1 & -2 & -1 \\ 0 & 0 & 0 \\ 1 & 2 & 1 \end{bmatrix}
\end{equation}
These kernels combine smoothing (perpendicular to gradient direction) with differentiation
(parallel to gradient direction), providing noise robustness.
\end{theorem}

\begin{remark}[Alternative Gradient Operators]
Other first-order operators include:
\begin{itemize}
\item \textbf{Prewitt}: Equal-weight smoothing, kernels $\begin{bmatrix} -1 & 0 & 1 \\ -1 & 0 & 1 \\ -1 & 0 & 1 \end{bmatrix}$
\item \textbf{Roberts}: 2$\times$2 diagonal kernels $\begin{bmatrix} 1 & 0 \\ 0 & -1 \end{bmatrix}$, $\begin{bmatrix} 0 & 1 \\ -1 & 0 \end{bmatrix}$
\item \textbf{Scharr}: Optimized 3$\times$3 kernels with coefficients $\pm 3, \pm 10$
\end{itemize}
\end{remark}

\subsection{Canny Edge Detector}

The Canny algorithm (1986) remains the gold standard for edge detection due to its optimal
performance according to three criteria: good detection, good localization, and single response.

\begin{theorem}[Canny Edge Detection Pipeline]
The Canny algorithm consists of four stages:
\begin{enumerate}
\item \textbf{Gaussian smoothing}: Convolve image with $G_\sigma(x,y) = \frac{1}{2\pi\sigma^2}e^{-(x^2+y^2)/2\sigma^2}$
\item \textbf{Gradient computation}: Calculate magnitude and orientation using Sobel or similar
\item \textbf{Non-maximum suppression}: Thin edges to single-pixel width by suppressing non-maxima perpendicular to edge direction
\item \textbf{Hysteresis thresholding}: Use dual thresholds $T_{low}$ and $T_{high}$ to link edge chains
\end{enumerate}
\end{theorem}

\subsection{Laplacian-Based Methods}

Second-order derivatives detect edges as zero-crossings rather than gradient maxima.

\begin{definition}[Laplacian of Gaussian (LoG)]
The LoG operator combines Gaussian smoothing with the Laplacian:
\begin{equation}
\nabla^2 G_\sigma = -\frac{1}{\pi\sigma^4}\left[1 - \frac{x^2+y^2}{2\sigma^2}\right]e^{-(x^2+y^2)/2\sigma^2}
\end{equation}
Edges occur where $\nabla^2(G_\sigma * I) = 0$ with sign change.
\end{definition}

\begin{remark}[Difference of Gaussians]
The LoG can be approximated by the Difference of Gaussians (DoG):
\begin{equation}
\text{DoG}(x,y) = G_{\sigma_1}(x,y) - G_{\sigma_2}(x,y), \quad \sigma_2 \approx 1.6\sigma_1
\end{equation}
This approximation is computationally efficient and forms the basis of SIFT features.
\end{remark}

\subsection{Performance Evaluation}

\begin{definition}[Edge Detection Metrics]
Given ground truth edges $E_{GT}$ and detected edges $E_{D}$:
\begin{itemize}
\item \textbf{Precision}: $P = \frac{TP}{TP + FP}$ (fraction of detected edges that are correct)
\item \textbf{Recall}: $R = \frac{TP}{TP + FN}$ (fraction of true edges detected)
\item \textbf{F-measure}: $F_1 = \frac{2PR}{P+R}$ (harmonic mean of precision and recall)
\item \textbf{Pratt's FOM}: $\text{FOM} = \frac{1}{\max(N_I, N_D)}\sum_{i=1}^{N_D}\frac{1}{1+\alpha d_i^2}$ where $d_i$ is distance to nearest true edge
\end{itemize}
\end{definition}

\section{Computational Analysis}

\begin{pycode}
import numpy as np
import matplotlib.pyplot as plt
from scipy import ndimage, signal
from scipy.ndimage import gaussian_filter, convolve

np.random.seed(42)

# Define gradient operators
sobel_x = np.array([[-1, 0, 1], [-2, 0, 2], [-1, 0, 1]])
sobel_y = np.array([[-1, -2, -1], [0, 0, 0], [1, 2, 1]])

prewitt_x = np.array([[-1, 0, 1], [-1, 0, 1], [-1, 0, 1]])
prewitt_y = np.array([[-1, -1, -1], [0, 0, 0], [1, 1, 1]])

roberts_x = np.array([[1, 0], [0, -1]])
roberts_y = np.array([[0, 1], [-1, 0]])

# Create synthetic test image: step edges at known locations
img_size = 200
test_image = np.zeros((img_size, img_size))
# Vertical edge
test_image[:, img_size//2:] = 1.0
# Horizontal edge in right half
test_image[img_size//2:, img_size//2:] = 0.5
# Add diagonal edge
for i in range(img_size//2):
    test_image[i, i:i+2] = 0.7

# Add Gaussian noise
noise_level = 0.05
noisy_image = test_image + noise_level * np.random.randn(img_size, img_size)
noisy_image = np.clip(noisy_image, 0, 1)

# Function implementations
def apply_sobel(image):
    gradient_x = convolve(image, sobel_x)
    gradient_y = convolve(image, sobel_y)
    gradient_magnitude = np.sqrt(gradient_x**2 + gradient_y**2)
    gradient_direction = np.arctan2(gradient_y, gradient_x)
    return gradient_magnitude, gradient_direction, gradient_x, gradient_y

def apply_prewitt(image):
    gradient_x = convolve(image, prewitt_x)
    gradient_y = convolve(image, prewitt_y)
    gradient_magnitude = np.sqrt(gradient_x**2 + gradient_y**2)
    return gradient_magnitude

def non_maximum_suppression(gradient_mag, gradient_dir):
    suppressed = np.zeros_like(gradient_mag)
    angle = gradient_dir * 180.0 / np.pi
    angle[angle < 0] += 180

    M, N = gradient_mag.shape
    for i in range(1, M-1):
        for j in range(1, N-1):
            q = 255
            r = 255

            # Angle 0 (horizontal)
            if (0 <= angle[i,j] < 22.5) or (157.5 <= angle[i,j] <= 180):
                q = gradient_mag[i, j+1]
                r = gradient_mag[i, j-1]
            # Angle 45
            elif (22.5 <= angle[i,j] < 67.5):
                q = gradient_mag[i+1, j-1]
                r = gradient_mag[i-1, j+1]
            # Angle 90 (vertical)
            elif (67.5 <= angle[i,j] < 112.5):
                q = gradient_mag[i+1, j]
                r = gradient_mag[i-1, j]
            # Angle 135
            elif (112.5 <= angle[i,j] < 157.5):
                q = gradient_mag[i-1, j-1]
                r = gradient_mag[i+1, j+1]

            if (gradient_mag[i,j] >= q) and (gradient_mag[i,j] >= r):
                suppressed[i,j] = gradient_mag[i,j]

    return suppressed

def hysteresis_thresholding(image, threshold_low, threshold_high):
    strong_edges = image > threshold_high
    weak_edges = (image >= threshold_low) & (image <= threshold_high)

    edges = strong_edges.copy()
    M, N = image.shape

    # Link weak edges to strong edges
    for i in range(1, M-1):
        for j in range(1, N-1):
            if weak_edges[i, j]:
                # Check 8-connected neighbors
                if np.any(strong_edges[i-1:i+2, j-1:j+2]):
                    edges[i, j] = True

    return edges.astype(float)

def canny_edge_detector(image, sigma=1.0, threshold_low=0.05, threshold_high=0.15):
    # Step 1: Gaussian smoothing
    smoothed = gaussian_filter(image, sigma)

    # Step 2: Gradient computation
    gradient_mag, gradient_dir, _, _ = apply_sobel(smoothed)

    # Step 3: Non-maximum suppression
    suppressed = non_maximum_suppression(gradient_mag, gradient_dir)

    # Step 4: Hysteresis thresholding
    edges = hysteresis_thresholding(suppressed, threshold_low, threshold_high)

    return edges, gradient_mag, suppressed

def laplacian_of_gaussian(image, sigma=2.0):
    # Smooth with Gaussian
    smoothed = gaussian_filter(image, sigma)
    # Apply Laplacian
    laplacian = ndimage.laplace(smoothed)
    return laplacian

def zero_crossing_detector(laplacian, threshold=0.01):
    zero_crossings = np.zeros_like(laplacian)
    M, N = laplacian.shape

    for i in range(1, M-1):
        for j in range(1, N-1):
            # Check for sign changes in 4-connected neighborhood
            region = laplacian[i-1:i+2, j-1:j+2]
            if (np.min(region) < -threshold) and (np.max(region) > threshold):
                zero_crossings[i, j] = 1

    return zero_crossings

# Apply different edge detectors
sobel_mag, sobel_dir, sobel_gx, sobel_gy = apply_sobel(noisy_image)
prewitt_mag = apply_prewitt(noisy_image)
canny_edges, canny_grad, canny_nms = canny_edge_detector(noisy_image, sigma=1.4,
                                                           threshold_low=0.05, threshold_high=0.15)
log_result = laplacian_of_gaussian(noisy_image, sigma=2.0)
log_edges = zero_crossing_detector(log_result, threshold=0.005)

# Threshold gradient magnitudes for comparison
sobel_edges = sobel_mag > 0.2
prewitt_edges = prewitt_mag > 0.2

# Count detected edges
num_sobel_edges = np.sum(sobel_edges)
num_prewitt_edges = np.sum(prewitt_edges)
num_canny_edges = np.sum(canny_edges)
num_log_edges = np.sum(log_edges)

# Calculate ground truth edges from clean image
true_edges_v = np.zeros_like(test_image)
true_edges_v[:, img_size//2-1:img_size//2+1] = 1
true_edges_h = np.zeros_like(test_image)
true_edges_h[img_size//2-1:img_size//2+1, img_size//2:] = 1
ground_truth = (true_edges_v + true_edges_h) > 0

# Compute performance metrics for Canny
def compute_metrics(detected, ground_truth, tolerance=2):
    # Dilate ground truth for tolerance
    from scipy.ndimage import binary_dilation
    struct = np.ones((2*tolerance+1, 2*tolerance+1))
    gt_dilated = binary_dilation(ground_truth, structure=struct)

    true_positive = np.sum(detected & gt_dilated)
    false_positive = np.sum(detected & ~gt_dilated)
    false_negative = np.sum(~detected & ground_truth)

    precision = true_positive / (true_positive + false_positive) if (true_positive + false_positive) > 0 else 0
    recall = true_positive / (true_positive + false_negative) if (true_positive + false_negative) > 0 else 0
    f_measure = 2 * precision * recall / (precision + recall) if (precision + recall) > 0 else 0

    return precision, recall, f_measure

canny_precision, canny_recall, canny_f1 = compute_metrics(canny_edges, ground_truth)
sobel_precision, sobel_recall, sobel_f1 = compute_metrics(sobel_edges, ground_truth)

# Create comprehensive visualization
fig = plt.figure(figsize=(16, 14))

# Plot 1: Original test image
ax1 = fig.add_subplot(4, 4, 1)
ax1.imshow(test_image, cmap='gray', interpolation='nearest')
ax1.set_title('Clean Test Image', fontsize=10)
ax1.axis('off')

# Plot 2: Noisy image
ax2 = fig.add_subplot(4, 4, 2)
ax2.imshow(noisy_image, cmap='gray', interpolation='nearest')
ax2.set_title(f'Noisy Image (SNR$\\approx${1/noise_level:.1f})', fontsize=10)
ax2.axis('off')

# Plot 3: Sobel gradient magnitude
ax3 = fig.add_subplot(4, 4, 3)
im3 = ax3.imshow(sobel_mag, cmap='hot', interpolation='nearest')
ax3.set_title('Sobel Gradient Magnitude', fontsize=10)
ax3.axis('off')
plt.colorbar(im3, ax=ax3, fraction=0.046)

# Plot 4: Gradient direction
ax4 = fig.add_subplot(4, 4, 4)
im4 = ax4.imshow(sobel_dir, cmap='hsv', interpolation='nearest')
ax4.set_title('Gradient Direction', fontsize=10)
ax4.axis('off')
plt.colorbar(im4, ax=ax4, fraction=0.046)

# Plot 5: Sobel X gradient
ax5 = fig.add_subplot(4, 4, 5)
im5 = ax5.imshow(sobel_gx, cmap='RdBu_r', interpolation='nearest')
ax5.set_title('Sobel $G_x$ (Vertical Edges)', fontsize=10)
ax5.axis('off')
plt.colorbar(im5, ax=ax5, fraction=0.046)

# Plot 6: Sobel Y gradient
ax6 = fig.add_subplot(4, 4, 6)
im6 = ax6.imshow(sobel_gy, cmap='RdBu_r', interpolation='nearest')
ax6.set_title('Sobel $G_y$ (Horizontal Edges)', fontsize=10)
ax6.axis('off')
plt.colorbar(im6, ax=ax6, fraction=0.046)

# Plot 7: Sobel thresholded
ax7 = fig.add_subplot(4, 4, 7)
ax7.imshow(sobel_edges, cmap='gray', interpolation='nearest')
ax7.set_title(f'Sobel Edges ({num_sobel_edges} pixels)', fontsize=10)
ax7.axis('off')

# Plot 8: Prewitt comparison
ax8 = fig.add_subplot(4, 4, 8)
ax8.imshow(prewitt_edges, cmap='gray', interpolation='nearest')
ax8.set_title(f'Prewitt Edges ({num_prewitt_edges} pixels)', fontsize=10)
ax8.axis('off')

# Plot 9: Canny - gradient after smoothing
ax9 = fig.add_subplot(4, 4, 9)
im9 = ax9.imshow(canny_grad, cmap='hot', interpolation='nearest')
ax9.set_title('Canny: Gradient (Smoothed)', fontsize=10)
ax9.axis('off')
plt.colorbar(im9, ax=ax9, fraction=0.046)

# Plot 10: Canny - after NMS
ax10 = fig.add_subplot(4, 4, 10)
im10 = ax10.imshow(canny_nms, cmap='hot', interpolation='nearest')
ax10.set_title('Canny: Non-Max Suppression', fontsize=10)
ax10.axis('off')
plt.colorbar(im10, ax=ax10, fraction=0.046)

# Plot 11: Canny final edges
ax11 = fig.add_subplot(4, 4, 11)
ax11.imshow(canny_edges, cmap='gray', interpolation='nearest')
ax11.set_title(f'Canny Edges ({num_canny_edges} pixels)', fontsize=10)
ax11.axis('off')

# Plot 12: Canny overlay
ax12 = fig.add_subplot(4, 4, 12)
overlay = np.stack([noisy_image, noisy_image, noisy_image], axis=-1)
overlay[canny_edges.astype(bool), :] = [1, 0, 0]
ax12.imshow(overlay, interpolation='nearest')
ax12.set_title('Canny Overlay on Image', fontsize=10)
ax12.axis('off')

# Plot 13: Laplacian of Gaussian
ax13 = fig.add_subplot(4, 4, 13)
im13 = ax13.imshow(log_result, cmap='RdBu_r', interpolation='nearest')
ax13.set_title('Laplacian of Gaussian', fontsize=10)
ax13.axis('off')
plt.colorbar(im13, ax=ax13, fraction=0.046)

# Plot 14: LoG zero crossings
ax14 = fig.add_subplot(4, 4, 14)
ax14.imshow(log_edges, cmap='gray', interpolation='nearest')
ax14.set_title(f'LoG Zero-Crossings ({num_log_edges} pixels)', fontsize=10)
ax14.axis('off')

# Plot 15: Performance comparison
ax15 = fig.add_subplot(4, 4, 15)
methods = ['Sobel', 'Canny']
precision_vals = [sobel_precision, canny_precision]
recall_vals = [sobel_recall, canny_recall]
f1_vals = [sobel_f1, canny_f1]

x_pos = np.arange(len(methods))
width = 0.25

ax15.bar(x_pos - width, precision_vals, width, label='Precision', color='steelblue', edgecolor='black')
ax15.bar(x_pos, recall_vals, width, label='Recall', color='coral', edgecolor='black')
ax15.bar(x_pos + width, f1_vals, width, label='F-measure', color='lightgreen', edgecolor='black')

ax15.set_ylabel('Score', fontsize=10)
ax15.set_title('Detection Performance', fontsize=10)
ax15.set_xticks(x_pos)
ax15.set_xticklabels(methods, fontsize=9)
ax15.legend(fontsize=8)
ax15.set_ylim([0, 1])
ax15.grid(axis='y', alpha=0.3)

# Plot 16: Edge count comparison
ax16 = fig.add_subplot(4, 4, 16)
edge_counts = [num_sobel_edges, num_prewitt_edges, num_canny_edges, num_log_edges]
count_labels = ['Sobel', 'Prewitt', 'Canny', 'LoG']
colors_bars = ['steelblue', 'coral', 'lightgreen', 'plum']

ax16.bar(count_labels, edge_counts, color=colors_bars, edgecolor='black')
ax16.set_ylabel('Number of Edge Pixels', fontsize=10)
ax16.set_title('Edge Detector Comparison', fontsize=10)
ax16.tick_params(axis='x', labelsize=9)
ax16.grid(axis='y', alpha=0.3)

plt.tight_layout()
plt.savefig('edge_detection_comprehensive.pdf', dpi=150, bbox_inches='tight')
plt.close()
\end{pycode}

\begin{figure}[htbp]
\centering
\includegraphics[width=\textwidth]{edge_detection_comprehensive.pdf}
\caption{Comprehensive edge detection analysis: (a) Clean synthetic test image with step edges;
(b) Image corrupted with Gaussian noise; (c-d) Sobel gradient magnitude and direction revealing
edge strength and orientation; (e-f) Directional gradients $G_x$ and $G_y$ highlighting vertical
and horizontal features; (g-h) Thresholded Sobel and Prewitt edges showing similar detection
patterns; (i-l) Canny detector pipeline stages from gradient computation through non-maximum
suppression to final hysteresis thresholding with overlay; (m-n) Laplacian of Gaussian response
and zero-crossing detection; (o-p) Quantitative performance comparison showing Canny's superior
precision-recall balance and total edge pixel counts across methods.}
\label{fig:edge_detection}
\end{figure}

\section{Algorithm Implementations}

\subsection{Canny Edge Detector}

\begin{algorithm}
\caption{Canny Edge Detection}
\begin{algorithmic}[1]
\Require Image $I$, Gaussian $\sigma$, thresholds $T_{low}$, $T_{high}$
\Ensure Binary edge map $E$
\State $I_{smooth} \gets G_\sigma * I$ \Comment{Gaussian smoothing}
\State $G_x \gets$ Sobel$_x * I_{smooth}$
\State $G_y \gets$ Sobel$_y * I_{smooth}$
\State $M \gets \sqrt{G_x^2 + G_y^2}$ \Comment{Gradient magnitude}
\State $\theta \gets \arctan(G_y / G_x)$ \Comment{Gradient direction}
\State $M_{nms} \gets$ NonMaxSuppress$(M, \theta)$ \Comment{Thin edges}
\State $E \gets$ HysteresisThreshold$(M_{nms}, T_{low}, T_{high})$ \Comment{Link edges}
\State \Return $E$
\end{algorithmic}
\end{algorithm}

\subsection{Non-Maximum Suppression}

The key to single-pixel-wide edges is suppressing gradient magnitudes that are not local maxima
in the direction perpendicular to the edge.

\begin{algorithm}
\caption{Non-Maximum Suppression}
\begin{algorithmic}[1]
\Require Gradient magnitude $M$, direction $\theta$
\Ensure Suppressed magnitude $M_{nms}$
\For{each pixel $(i,j)$}
    \State Quantize $\theta_{i,j}$ to nearest 45$^\circ$ (0, 45, 90, 135)
    \State Get neighboring pixels $p, q$ along gradient direction
    \If{$M_{i,j} \geq M_p$ \textbf{and} $M_{i,j} \geq M_q$}
        \State $M_{nms}[i,j] \gets M_{i,j}$ \Comment{Local maximum}
    \Else
        \State $M_{nms}[i,j] \gets 0$ \Comment{Suppressed}
    \EndIf
\EndFor
\State \Return $M_{nms}$
\end{algorithmic}
\end{algorithm}

\section{Natural Image Analysis}

\begin{pycode}
# Create a more complex natural-like scene
np.random.seed(123)
scene_size = 300

# Create synthetic natural scene with multiple objects
scene = np.ones((scene_size, scene_size)) * 0.3

# Add geometric shapes with different intensities
# Circle
center_y, center_x = 100, 100
radius = 40
y, x = np.ogrid[:scene_size, :scene_size]
circle_mask = (x - center_x)**2 + (y - center_y)**2 <= radius**2
scene[circle_mask] = 0.8

# Rectangle
scene[150:220, 50:150] = 0.6

# Triangle (approximate)
for i in range(70):
    scene[200-i, 180:180+2*i] = 0.45

# Add texture using random noise patterns
texture = np.random.randn(scene_size, scene_size) * 0.03
scene += texture
scene = np.clip(scene, 0, 1)

# Add realistic noise
scene_noisy = scene + 0.02 * np.random.randn(scene_size, scene_size)
scene_noisy = np.clip(scene_noisy, 0, 1)

# Apply multiple edge detectors with different parameters
canny_sigma1, _, _ = canny_edge_detector(scene_noisy, sigma=0.8, threshold_low=0.04, threshold_high=0.12)
canny_sigma2, _, _ = canny_edge_detector(scene_noisy, sigma=1.5, threshold_low=0.04, threshold_high=0.12)
canny_sigma3, _, _ = canny_edge_detector(scene_noisy, sigma=2.5, threshold_low=0.04, threshold_high=0.12)

sobel_scene, _, _, _ = apply_sobel(scene_noisy)
log_scene = laplacian_of_gaussian(scene_noisy, sigma=1.5)
log_scene_edges = zero_crossing_detector(log_scene, threshold=0.003)

# Multi-scale analysis
scales = [0.5, 1.0, 2.0, 3.0]
multiscale_edges = np.zeros_like(scene)

for sigma in scales:
    edges, _, _ = canny_edge_detector(scene_noisy, sigma=sigma, threshold_low=0.04, threshold_high=0.12)
    multiscale_edges += edges

multiscale_edges = multiscale_edges > 0

# Create visualization
fig2 = plt.figure(figsize=(16, 10))

# Original scene
ax1 = fig2.add_subplot(3, 4, 1)
ax1.imshow(scene, cmap='gray', interpolation='nearest')
ax1.set_title('Synthetic Scene (Clean)', fontsize=10)
ax1.axis('off')

# Noisy scene
ax2 = fig2.add_subplot(3, 4, 2)
ax2.imshow(scene_noisy, cmap='gray', interpolation='nearest')
ax2.set_title('Noisy Scene', fontsize=10)
ax2.axis('off')

# Sobel magnitude
ax3 = fig2.add_subplot(3, 4, 3)
im3 = ax3.imshow(sobel_scene, cmap='hot', interpolation='nearest')
ax3.set_title('Sobel Gradient Magnitude', fontsize=10)
ax3.axis('off')
plt.colorbar(im3, ax=ax3, fraction=0.046)

# Canny sigma=0.8
ax4 = fig2.add_subplot(3, 4, 4)
ax4.imshow(canny_sigma1, cmap='gray', interpolation='nearest')
ax4.set_title(f'Canny ($\\sigma$=0.8, {np.sum(canny_sigma1)} edges)', fontsize=10)
ax4.axis('off')

# Canny sigma=1.5
ax5 = fig2.add_subplot(3, 4, 5)
ax5.imshow(canny_sigma2, cmap='gray', interpolation='nearest')
ax5.set_title(f'Canny ($\\sigma$=1.5, {np.sum(canny_sigma2)} edges)', fontsize=10)
ax5.axis('off')

# Canny sigma=2.5
ax6 = fig2.add_subplot(3, 4, 6)
ax6.imshow(canny_sigma3, cmap='gray', interpolation='nearest')
ax6.set_title(f'Canny ($\\sigma$=2.5, {np.sum(canny_sigma3)} edges)', fontsize=10)
ax6.axis('off')

# LoG response
ax7 = fig2.add_subplot(3, 4, 7)
im7 = ax7.imshow(log_scene, cmap='RdBu_r', interpolation='nearest')
ax7.set_title('LoG Response ($\\sigma$=1.5)', fontsize=10)
ax7.axis('off')
plt.colorbar(im7, ax=ax7, fraction=0.046)

# LoG edges
ax8 = fig2.add_subplot(3, 4, 8)
ax8.imshow(log_scene_edges, cmap='gray', interpolation='nearest')
ax8.set_title(f'LoG Edges ({np.sum(log_scene_edges)} pixels)', fontsize=10)
ax8.axis('off')

# Multi-scale edges
ax9 = fig2.add_subplot(3, 4, 9)
ax9.imshow(multiscale_edges, cmap='gray', interpolation='nearest')
ax9.set_title('Multi-Scale Edges (Combined)', fontsize=10)
ax9.axis('off')

# Edge count vs sigma
ax10 = fig2.add_subplot(3, 4, 10)
sigma_range = np.linspace(0.5, 4.0, 15)
edge_counts_sigma = []
for sig in sigma_range:
    edges_temp, _, _ = canny_edge_detector(scene_noisy, sigma=sig, threshold_low=0.04, threshold_high=0.12)
    edge_counts_sigma.append(np.sum(edges_temp))

ax10.plot(sigma_range, edge_counts_sigma, 'o-', linewidth=2, markersize=6, color='steelblue')
ax10.set_xlabel('Gaussian $\\sigma$', fontsize=10)
ax10.set_ylabel('Edge Pixels Detected', fontsize=10)
ax10.set_title('Scale Parameter Effect', fontsize=10)
ax10.grid(True, alpha=0.3)

# Threshold sensitivity
ax11 = fig2.add_subplot(3, 4, 11)
threshold_range = np.linspace(0.02, 0.20, 15)
edge_counts_thresh = []
for thresh in threshold_range:
    edges_temp, _, _ = canny_edge_detector(scene_noisy, sigma=1.5,
                                            threshold_low=thresh, threshold_high=thresh*3)
    edge_counts_thresh.append(np.sum(edges_temp))

ax11.plot(threshold_range, edge_counts_thresh, 's-', linewidth=2, markersize=6, color='coral')
ax11.set_xlabel('Threshold $T_{low}$', fontsize=10)
ax11.set_ylabel('Edge Pixels Detected', fontsize=10)
ax11.set_title('Threshold Sensitivity', fontsize=10)
ax11.grid(True, alpha=0.3)

# Histogram of gradient magnitudes
ax12 = fig2.add_subplot(3, 4, 12)
gradient_flat = sobel_scene.flatten()
ax12.hist(gradient_flat, bins=50, density=True, alpha=0.7, color='steelblue', edgecolor='black')
ax12.axvline(x=0.12, color='red', linestyle='--', linewidth=2, label='Typical Threshold')
ax12.set_xlabel('Gradient Magnitude', fontsize=10)
ax12.set_ylabel('Probability Density', fontsize=10)
ax12.set_title('Gradient Distribution', fontsize=10)
ax12.legend(fontsize=9)
ax12.set_xlim([0, 0.5])

plt.tight_layout()
plt.savefig('edge_detection_natural_scene.pdf', dpi=150, bbox_inches='tight')
plt.close()

# Store key statistics
num_canny_sigma1 = int(np.sum(canny_sigma1))
num_canny_sigma2 = int(np.sum(canny_sigma2))
num_canny_sigma3 = int(np.sum(canny_sigma3))
num_log_scene = int(np.sum(log_scene_edges))
num_multiscale = int(np.sum(multiscale_edges))
\end{pycode}

\begin{figure}[htbp]
\centering
\includegraphics[width=\textwidth]{edge_detection_natural_scene.pdf}
\caption{Edge detection on synthetic natural scene: (a-b) Clean and noisy synthetic scenes
containing geometric shapes with varying intensities; (c) Sobel gradient magnitude showing
continuous edge strength; (d-f) Canny edge detection at three scales ($\sigma$ = 0.8, 1.5, 2.5)
demonstrating the noise-detail trade-off, detecting \py{num_canny_sigma1}, \py{num_canny_sigma2},
and \py{num_canny_sigma3} edge pixels respectively; (g-h) Laplacian of Gaussian response and
zero-crossing detection finding \py{num_log_scene} edge pixels; (i) Multi-scale edge combination
yielding \py{num_multiscale} consolidated edge pixels; (j-k) Parameter sensitivity analysis showing
edge count dependence on Gaussian smoothing scale and threshold values; (l) Gradient magnitude
histogram revealing the distribution of edge strengths and typical threshold selection criterion.}
\label{fig:natural_scene}
\end{figure}

\section{Results and Performance Analysis}

\subsection{Computational Complexity}

\begin{pycode}
print(r"\begin{table}[htbp]")
print(r"\centering")
print(r"\caption{Edge Detection Algorithm Complexity}")
print(r"\begin{tabular}{lccc}")
print(r"\toprule")
print(r"Algorithm & Convolutions & Non-linear Ops & Complexity \\")
print(r"\midrule")
print(r"Sobel & 2 ($3\times3$) & Sqrt, Atan & $O(n^2)$ \\")
print(r"Prewitt & 2 ($3\times3$) & Sqrt, Atan & $O(n^2)$ \\")
print(r"Roberts & 2 ($2\times2$) & Sqrt, Atan & $O(n^2)$ \\")
print(r"Canny & 3 (Gaussian + 2 Sobel) & NMS + Hysteresis & $O(n^2)$ \\")
print(r"LoG & 1 (LoG kernel) & Zero-crossing & $O(n^2)$ \\")
print(r"\bottomrule")
print(r"\end{tabular}")
print(r"\label{tab:complexity}")
print(r"\end{table}")
\end{pycode}

\subsection{Detection Performance Summary}

\begin{pycode}
print(r"\begin{table}[htbp]")
print(r"\centering")
print(r"\caption{Edge Detector Performance on Test Image}")
print(r"\begin{tabular}{lcccc}")
print(r"\toprule")
print(r"Method & Edges Detected & Precision & Recall & F-measure \\")
print(r"\midrule")
print(f"Sobel (threshold=0.2) & {num_sobel_edges} & {sobel_precision:.3f} & {sobel_recall:.3f} & {sobel_f1:.3f} \\\\")
print(f"Prewitt (threshold=0.2) & {num_prewitt_edges} & --- & --- & --- \\\\")
print(f"Canny ($\\sigma$=1.4) & {num_canny_edges} & {canny_precision:.3f} & {canny_recall:.3f} & {canny_f1:.3f} \\\\")
print(f"LoG ($\\sigma$=2.0) & {num_log_edges} & --- & --- & --- \\\\")
print(r"\bottomrule")
print(r"\end{tabular}")
print(r"\label{tab:performance}")
print(r"\end{table}")
\end{pycode}

\subsection{Scale Parameter Analysis}

The choice of Gaussian smoothing scale $\sigma$ represents a fundamental trade-off in edge detection:

\begin{itemize}
\item \textbf{Small $\sigma$ (0.5-1.0)}: Preserves fine details but sensitive to noise, detects $\sim$\py{num_canny_sigma1} edges
\item \textbf{Medium $\sigma$ (1.5-2.5)}: Balanced noise suppression and detail retention, detects $\sim$\py{num_canny_sigma2} edges
\item \textbf{Large $\sigma$ (3.0-4.0)}: Strong noise rejection but loses fine structure, detects $\sim$\py{num_canny_sigma3} edges
\end{itemize}

\section{Discussion}

\begin{example}[Gradient Operator Selection]
For real-time applications on noisy images, the Sobel operator provides an excellent balance:
\begin{itemize}
\item \textbf{Efficiency}: Only two $3\times3$ convolutions required
\item \textbf{Robustness}: Built-in smoothing suppresses high-frequency noise
\item \textbf{Accuracy}: Weighted gradient approximation superior to simple differences
\end{itemize}
For highest quality, Canny's multi-stage pipeline achieves F-measure of \py{f"{canny_f1:.3f}"}
compared to Sobel's \py{f"{sobel_f1:.3f}"} on our test image.
\end{example}

\begin{remark}[Second vs First Derivatives]
Laplacian-based methods offer theoretical advantages:
\begin{itemize}
\item Isotropic response (rotation invariant)
\item Precise localization at zero-crossings
\item No explicit thresholding of gradient magnitude
\end{itemize}
However, second derivatives amplify noise more severely, requiring larger smoothing kernels
and potentially missing low-contrast edges.
\end{remark}

\begin{example}[Hysteresis Thresholding]
Canny's dual-threshold approach solves the edge linking problem:
\begin{itemize}
\item High threshold ($T_h$): Seeds strong edge segments (high confidence)
\item Low threshold ($T_l$): Extends edges along weak gradients (connected to strong edges)
\item Typical ratio: $T_h/T_l \approx 2-3$
\end{itemize}
This prevents edge fragmentation while rejecting isolated noise pixels.
\end{example}

\section{Conclusions}

This analysis demonstrates the implementation and evaluation of fundamental edge detection algorithms:

\begin{enumerate}
\item Gradient-based methods (Sobel, Prewitt) detect \py{num_sobel_edges} and \py{num_prewitt_edges}
      edge pixels respectively on the noisy test image
\item The Canny detector's multi-stage pipeline achieves superior performance with precision
      \py{f"{canny_precision:.3f}"}, recall \py{f"{canny_recall:.3f}"}, and F-measure \py{f"{canny_f1:.3f}"}
\item Laplacian of Gaussian methods provide rotation-invariant detection with \py{num_log_edges}
      zero-crossings identified
\item Scale selection critically impacts detection: smaller $\sigma$ values preserve detail
      (\py{num_canny_sigma1} edges at $\sigma$=0.8) while larger values enhance noise robustness
      (\py{num_canny_sigma3} edges at $\sigma$=2.5)
\item Multi-scale approaches combining detections across $\sigma \in [0.5, 3.0]$ yield
      \py{num_multiscale} consolidated edge pixels
\end{enumerate}

The choice of edge detector depends on application requirements: Sobel for computational efficiency,
Canny for accuracy and completeness, and LoG for precise localization in low-noise scenarios.

\section*{References}

\begin{enumerate}
\item Canny, J. (1986). A computational approach to edge detection. \textit{IEEE Trans. Pattern Analysis and Machine Intelligence}, 8(6), 679-698.
\item Sobel, I., \& Feldman, G. (1968). A 3x3 isotropic gradient operator for image processing. \textit{Stanford Artificial Intelligence Project}.
\item Marr, D., \& Hildreth, E. (1980). Theory of edge detection. \textit{Proc. Royal Society of London B}, 207(1167), 187-217.
\item Prewitt, J. M. S. (1970). Object enhancement and extraction. \textit{Picture Processing and Psychopictorics}, Academic Press.
\item Roberts, L. G. (1965). Machine perception of three-dimensional solids. \textit{Optical and Electro-Optical Information Processing}, MIT Press.
\item Torre, V., \& Poggio, T. A. (1986). On edge detection. \textit{IEEE Trans. Pattern Analysis and Machine Intelligence}, 8(2), 147-163.
\item Haralick, R. M. (1984). Digital step edges from zero crossing of second directional derivatives. \textit{IEEE Trans. Pattern Analysis and Machine Intelligence}, 6(1), 58-68.
\item Scharr, H. (2000). Optimal operators in digital image processing. \textit{PhD Thesis}, University of Heidelberg.
\item Elder, J. H., \& Zucker, S. W. (1998). Local scale control for edge detection and blur estimation. \textit{IEEE Trans. Pattern Analysis and Machine Intelligence}, 20(7), 699-716.
\item Lindeberg, T. (1998). Edge detection and ridge detection with automatic scale selection. \textit{Int. Journal of Computer Vision}, 30(2), 117-156.
\item Ziou, D., \& Tabbone, S. (1998). Edge detection techniques: An overview. \textit{International Journal of Pattern Recognition and Image Analysis}, 8(4), 537-559.
\item Heath, M., et al. (1997). A robust visual method for assessing the relative performance of edge-detection algorithms. \textit{IEEE Trans. Pattern Analysis and Machine Intelligence}, 19(12), 1338-1359.
\item Kittler, J. (1983). On the accuracy of the Sobel edge detector. \textit{Image and Vision Computing}, 1(1), 37-42.
\item Deriche, R. (1987). Using Canny's criteria to derive a recursively implemented optimal edge detector. \textit{Int. Journal of Computer Vision}, 1(2), 167-187.
\item Bovik, A. C., Huang, T. S., \& Munson, D. C. (1983). A generalization of median filtering using linear combinations of order statistics. \textit{IEEE Trans. Acoustics, Speech, and Signal Processing}, 31(6), 1342-1350.
\item Pratt, W. K. (2007). \textit{Digital Image Processing}, 4th ed. Wiley-Interscience.
\item Gonzalez, R. C., \& Woods, R. E. (2018). \textit{Digital Image Processing}, 4th ed. Pearson.
\item Jain, A. K. (1989). \textit{Fundamentals of Digital Image Processing}. Prentice Hall.
\end{enumerate}

\end{document}
